% GENERATED FILE. Edit canonical lessons, then run npm run generate:books.
% canonical-source-hash: fnv1a64:5aae9a95cf256abe
% canonical-lessons: ZH-C353-shuofu, ZH-C353-daoli, ZH-W353-dao-way, ZH-C353-bujiande, ZH-C353-fanbo, ZH-C353-rentong

\chapter{You've Got a Point}
\label{ch:zh-you-ve-got-a-point}
\hlchaptermodality{353}

\begin{chapteropening}
\textbf{By the end of this chapter:} \emph{I can say someone has a point, talk someone round and say something is not necessarily so.}

Complete the last lesson of chapter 353: I can say someone has a point, talk someone round and say something is not necessarily so.
\end{chapteropening}

\section[shuōfú]{shuōfú — to talk someone round; to persuade}

\label{lesson:ZH-C353-shuofu}



\begin{quote}
Before the new one: say the Mandarin for what one really meant, then the Mandarin for to make your view known.
\end{quote}

\subsection*{You'll want to know: shuōfú}
\emph{shuōfú} — \textquotedblleft{}to talk someone round; to persuade\textquotedblright{}.

\begin{quote}
\emph{wǒ juéde nǐ néng shuōfú tā} — I think you can talk him round.
\end{quote}

Its characters come later; for now it is a word for the ear.

\begin{cousinweb}
\emph{shuōfú} is built from \zh{说} 'to say' and \emph{fú}, both inherited from Middle Chinese, \emph{sywet} and \emph{bjuwk} (Guangyun, in Baxter's transcription), which glosses the first 'to tell' and the second 'to serve' and 'clothes'. How the pair came to mean persuading is not settled.
\end{cousinweb}

\begin{grammarlens}[title={What you've built}]
One more word for giving reasons.
\end{grammarlens}

\subsection*{Your turn}
\begin{itemize}
\raggedright
  \item \emph{Say it:} \emph{shuōfú}
  \item \emph{Say it:} \emph{shuōfú}, once more
  \item \emph{Say it:} \emph{wǒ juéde nǐ néng shuōfú tā} — I think you can talk him round.
  \item \emph{Recall:} say the Mandarin for what one really meant, then the Mandarin for to make your view known, then say \emph{shuōfú} again
\end{itemize}

\begin{tcolorbox}[breakable,colback=teal!4,colframe=teal!35!black,title={Before you move on}]
What does shuōfú mean? (\textquotedblleft{}To talk someone round; to persuade\textquotedblright{}.) Say it once more.
\end{tcolorbox}
\section[\texorpdfstring{\zh{道理} (dàoli)}{道理 (dàoli)}]{\zh{道理} (dàoli) — sense; the reason something makes sense}

\label{lesson:ZH-C353-daoli}



\begin{quote}
Before the new one: say the Mandarin for to make your view known, then the Mandarin for to talk someone round.
\end{quote}

\subsection*{You'll want to know: \zh{道理}}
\textbf{\zh{道理}} — \emph{dàoli} — \textquotedblleft{}sense; the reason something makes sense\textquotedblright{}.

\begin{quote}
\textbf{\zh{你说得很有道理}} \emph{nǐ shuō de hěn yǒu dàoli} — You've got a point there.
\end{quote}

The new character, in pieces: \zh{道} — twelve strokes in two pieces: first a head piece of nine, two short strokes, a bar, a short left-falling stroke and a narrow box with two bars inside; then the walking piece of \zh{这} and \zh{过}, three strokes, written last. The next lesson puts it in your hand.

\begin{cousinweb}
\textbf{\zh{道理}} (\emph{dàoli}) is built from \zh{道} 'a way' and \zh{理} 'order', both inherited from Middle Chinese, \emph{dawX} and \emph{liX} (Guangyun, in Baxter's transcription), which glosses the first 'a principle' and 'a road', and the second 'to put in order'. The second is said light here.
\end{cousinweb}

\begin{grammarlens}[title={What you've built}]
One more word for giving reasons.
\end{grammarlens}

\subsection*{Your turn}
\begin{itemize}
\raggedright
  \item \emph{Say it:} \emph{dàoli}
  \item \emph{Say it:} \emph{dàoli}, once more
  \item \emph{Say it:} \emph{nǐ shuō de hěn yǒu dàoli} — You've got a point there.
  \item \emph{Recall:} say the Mandarin for to make your view known, then the Mandarin for to talk someone round, then say \emph{dàoli} again
\end{itemize}

\begin{tcolorbox}[breakable,colback=teal!4,colframe=teal!35!black,title={Before you move on}]
What does \zh{道理} mean? (\textquotedblleft{}Sense; the reason something makes sense\textquotedblright{}.) Say it once more.
\end{tcolorbox}
\section[\texorpdfstring{\zh{道} (dào)}{道 (dào)}]{\zh{道} — a head piece, then the walking piece}

\label{lesson:ZH-W353-dao-way}



\begin{quote}
The word first, then the character inside it.

\begin{itemize}
\raggedright
  \item \emph{Recall:} say \emph{dàoli} — \textbf{R1}, one lesson back, before it decays
\end{itemize}

\begin{quote}
\zh{这} \zh{过} \zh{理}
\end{quote}

\emph{Read it:} each one aloud, with nothing said first — \emph{zhè}, \emph{guò}, \emph{lǐ}
\end{quote}

\subsection*{Script — a head piece, then the walking piece}
\hlblockfigure{\detokenize{figures/ZH-W353-dao-way-filmstrip.pdf}}{How \zh{道} is written, stroke by stroke}

\begin{quote}
\zh{道}
\end{quote}

\textbf{Two pieces: the head piece first, the walking piece last}, although it sits on the left and along the foot. Count the pieces, not the strokes.

\textbf{The head piece, nine strokes.} A short stroke falling to the right and a short stroke falling to the left, like two horns. A horizontal under them. A short left-falling stroke. Then a narrow box: a vertical, a horizontal that turns down, two short horizontals inside and one closing the foot.

\textbf{The walking piece, three strokes}, as in \zh{这}: a dot, a stroke that zigzags down, then a long sweep along the foot to the right, under the head piece.

9 + 3 = twelve.

\subsection*{Your turn}
\begin{itemize}
\raggedright
  \item \emph{Write it:} \zh{道}, all twelve, with the model in front of you
  \item \emph{Cover:} write it again from memory, and count the pieces, not the strokes
\end{itemize}

\begin{tcolorbox}[breakable,colback=teal!4,colframe=teal!35!black,title={Before you move on}]
How many pieces in \zh{道}? (\textbf{Two}.) Which is written first? (\textbf{The head piece}, though the walking piece is on the left.) Where have you written the walking piece? (\textbf{In \zh{这} and \zh{过}}.)

Source: \href{https://raw.githubusercontent.com/chanind/hanzi-writer-data/68d10a4b21150cae5e1ebbd223eed289cf32d90c/data/\%E9\%81\%93.json}{Hanzi Writer Data \zh{道}, pinned stroke paths and medians}.
\end{tcolorbox}
\section[\texorpdfstring{\zh{不见得} (bújiànde)}{不见得 (bújiànde)}]{\zh{不见得} (bújiànde) — not necessarily}

\label{lesson:ZH-C353-bujiande}



\begin{quote}
Before the new one: say the Mandarin for to talk someone round, then the Mandarin for the sense something makes.
\end{quote}

\subsection*{You'll want to know: \zh{不见得}}
\textbf{\zh{不见得}} — \emph{bújiànde} — \textquotedblleft{}not necessarily\textquotedblright{}.

\begin{quote}
\emph{guì de dōngxi bújiànde hǎo} — Expensive things aren't necessarily good.
\end{quote}

The tone of bù changes here: before a falling tone it is said bú, and the spelling writes the change.

\begin{cousinweb}
\textbf{\zh{不见得}} (\emph{bújiànde}) is built from \zh{不} 'not', \zh{见} 'to see' and \zh{得} 'to get', all three inherited from Middle Chinese, \emph{pjut}, \emph{kenH} and \emph{tok} (Guangyun, in Baxter's transcription), which glosses the first as another word for not, the second 'to see' and the third 'gain'. How the three came to mean not necessarily is not settled.
\end{cousinweb}

\begin{grammarlens}[title={What you've built}]
One more word for giving reasons.
\end{grammarlens}

\subsection*{Your turn}
\begin{itemize}
\raggedright
  \item \emph{Say it:} \emph{bújiànde}
  \item \emph{Say it:} \emph{bújiànde}, once more
  \item \emph{Say it:} \emph{guì de dōngxi bújiànde hǎo} — Expensive things aren't necessarily good.
  \item \emph{Recall:} say the Mandarin for to talk someone round, then the Mandarin for the sense something makes, then say \emph{bújiànde} again
\end{itemize}

\begin{tcolorbox}[breakable,colback=teal!4,colframe=teal!35!black,title={Before you move on}]
What does \zh{不见得} mean? (\textquotedblleft{}Not necessarily\textquotedblright{}.) Say it once more.
\end{tcolorbox}
\section[fǎnbó]{fǎnbó — to argue back; to refute}

\label{lesson:ZH-C353-fanbo}



\begin{quote}
Before the new one: say the Mandarin for the sense something makes, then the Mandarin for not necessarily.
\end{quote}

\subsection*{You'll want to know: fǎnbó}
\emph{fǎnbó} — \textquotedblleft{}to argue back; to refute\textquotedblright{}.

\begin{quote}
\emph{yīnwèi tā shuō de hěn yǒu dàoli suǒyǐ wǒ méiyǒu fǎnbó} — He had a point, so I didn't argue back.
\end{quote}

Its characters come later; for now it is a word for the ear.

\begin{cousinweb}
\emph{fǎnbó} is built from \emph{fǎn} and \emph{bó}, both inherited from Middle Chinese, \emph{pjonX} and \emph{paewk} (Guangyun, in Baxter's transcription), which glosses the first 'to turn over' and the second 'a horse of mixed colour'. How the pair came to mean arguing back is not settled.
\end{cousinweb}

\begin{grammarlens}[title={What you've built}]
One more word for giving reasons.
\end{grammarlens}

\subsection*{Your turn}
\begin{itemize}
\raggedright
  \item \emph{Say it:} \emph{fǎnbó}
  \item \emph{Say it:} \emph{fǎnbó}, once more
  \item \emph{Say it:} \emph{yīnwèi tā shuō de hěn yǒu dàoli suǒyǐ wǒ méiyǒu fǎnbó} — He had a point, so I didn't argue back.
  \item \emph{Recall:} say the Mandarin for the sense something makes, then the Mandarin for not necessarily, then say \emph{fǎnbó} again
\end{itemize}

\begin{tcolorbox}[breakable,colback=teal!4,colframe=teal!35!black,title={Before you move on}]
What does fǎnbó mean? (\textquotedblleft{}To argue back; to refute\textquotedblright{}.) Say it once more.
\end{tcolorbox}
\section[rèntóng]{rèntóng — to share someone's view; to identify with}

\label{lesson:ZH-C353-rentong}



\begin{quote}
Before the new one: say the Mandarin for not necessarily, then the Mandarin for to argue back.
\end{quote}

\subsection*{You'll want to know: rèntóng}
\emph{rèntóng} — \textquotedblleft{}to share someone's view; to identify with\textquotedblright{}.

\begin{quote}
\emph{wǒ hěn rèntóng nǐ de kànfǎ} — I really share your view.
\end{quote}

Its characters come later; for now it is a word for the ear.

\begin{cousinweb}
\emph{rèntóng} is built from \emph{rèn} and \zh{同} 'the same', both inherited from Middle Chinese, \emph{nyinH} and \emph{duwng} (Guangyun, in Baxter's transcription), which glosses the first 'to recognize' and the second 'equal'.
\end{cousinweb}

\begin{grammarlens}[title={What you've built}]
One more word for giving reasons.
\end{grammarlens}

\subsection*{Your turn}
\begin{itemize}
\raggedright
  \item \emph{Say it:} \emph{rèntóng}
  \item \emph{Say it:} \emph{rèntóng}, once more
  \item \emph{Say it:} \emph{wǒ hěn rèntóng nǐ de kànfǎ} — I really share your view.
  \item \emph{Recall:} say the Mandarin for not necessarily, then the Mandarin for to argue back, then say \emph{rèntóng} again
\end{itemize}

\begin{tcolorbox}[breakable,colback=teal!4,colframe=teal!35!black,title={Before you move on}]
What does rèntóng mean? (\textquotedblleft{}To share someone's view; to identify with\textquotedblright{}.) Say it once more.
\end{tcolorbox}
